\chapter{为什么造出一组零，就已经快结束了}\label{sec:blocks}
\section{先解决十五元五平均，看看矩阵从哪里来}
设十五个位置已是 $(a^5,b^5,c^5)$，且 $a+b+c=0$。按照下面的表重新分成三组：
\[
K_+=\begin{pmatrix}0&3&2\\1&2&2\\4&0&1\end{pmatrix}.
\]
第一行说“从 $a$ 块拿零份，从 $b$ 块拿三份，从 $c$ 块拿两份”；每行和为五，所以每行可以做一次五平均。每列和也为五，所以恰好用完每个旧块，没有多拿、漏拿或复制位置。

三个新值就是 $K_+(a,b,c)^T/5$。令 $x=a,y=a-b$，则 $b=x-y,c=y-2x$。代入第一行，得到
\[a'=-\frac{x+y}{5};
\]
第二行给 $b'=-x/5$，所以 $y'=a'-b'=-y/5$。这三次真实平均，在两个参数上就是
\[(x,y)\longmapsto-\frac15(x+y,y).
\]
暂时忽略整体比例，它把第一坐标加上第二坐标，这叫\textbf{上剪切}。

另一张表
\[K_-=\begin{pmatrix}2&1&2\\3&0&2\\0&4&1\end{pmatrix}
\]
实现 $-(x-y,y)/5$。两个词接起来，得到 $(x,y)/25$。参数的比例回来了，但真实数值整体缩小了；因此“方向上的逆”完全可以由正向平均实现。

\section{一般的三个等大块}
一般 $p\ge5$ 时写 $p=3m+\epsilon$、$\epsilon=\pm1$。在零和空间上，全一矩阵 $J$ 的作用为零。于是我们可以先设计想要的作用 $I\pm D$，再加 $mJ$，把所有计数变成非负。加上去的部分对零和参数没有影响，却让抽象矩阵成为真的分组。

下面会出现 $\SL_2(\Z)$，它只是所有行列式为一的二阶整数矩阵。$U(k)$ 把 $(x,y)$ 变为 $(x+ky,y)$；$L(k)$ 把它变为 $(x,y+kx)$。欧几里得算法正是在反复做这种整数加减。

$\Gamma_0(3,\Z)$ 是这些矩阵中左下项被三整除的那些。证明中的四陪集表，你可以看成一张四个状态的路线表：读入一个生成元后，记录落在哪个状态，并把差额记成已有的 $U,L$。读完整个矩阵词，就得到可以实际执行的剪切词。
\InsertProof{prop:threep}

\section{一个零给我们留了什么余地}
在 $2p+1$ 个零和位置中，如果有一个零，把它留着，其他 $2p$ 个位置正好分两组平均。两个组和相反，所以得到 $(a^p,(-a)^p,0)$。

例如 $p=5$，先拿两份 $a$、两份 $-a$、一个零平均，产生五个零；再做一次，只剩一对 $a,-a$；最后拿这对和三个零。只需三次就消去两个五重块。
\InsertProof{lem:zero-base}

\section{为什么可以在较小子系统里收尾}
一个零和 $p$ 元组平均以后，至少出现 $p$ 个零，所以非零位置最多有 $n-p$ 个。如果我们选一个子系统，\emph{包含全部非零位置，再多带一个零}，它的总和仍然是零，目标与整个配置一致。

含零还带来一个算术好处：所有坐标差的最大公因数等于坐标本身的最大公因数。因为差里面就包含 $x_i-0=x_i$。本原化后，这个子系统的 $G$ 必为一。

因此，只要这个子系统的大小属于已经证明的范围，就能在里面完成，外面的零保持不动。下面按 $n$ 的范围选择 $2p+1,3p,4p$ 等已解决的大小。
\InsertProof{lem:zero-tail}

后面每次说“送到终端”，具体目标都是先让某个 $p$ 元组的和恰好为零。其余操作由本讲接手。这样我们不必要求一个可逆矩阵把非零列直接变成零；它也确实做不到。
